\subsection{Shared experience banks and dependency-aware skill retrieval narrow the residual further}

A closer parent for the ``learn once, reuse across tasks'' intuition is ReX, which treats recurring reasoning patterns and skills as latent experiences rather than task-specific adapters \cite{ling2026rex}. ReX learns a shared Experience Bank of foundational skill vectors, encodes each input into a low-dimensional experience code, and dynamically composes the relevant skill vectors into a lightweight adapter without requiring explicit task identifiers. Its multi-task results show that cross-input/task experience reuse can improve specialization and generalization. This removes any novelty claim that Paper 3 uniquely proposes a persistent shared experience substrate or dynamic composition of reusable skills.

Dependency-aware skill retrieval is also moving beyond flat semantic matching. SkillGraph represents agent skills as a dependency graph and targets selection/composition failures that arise when modular skills are treated as isolated text descriptions \cite{wu2026skillgraph}. AgentGL similarly treats graph topology as an active reasoning substrate: an LLM agent uses graph-native tools at inference while a graph-conditioned curriculum is used during reinforcement learning \cite{sun2026agentgl}. These systems are important parent mechanisms because they show that relational topology can inform both learning and execution.

The candidate residual must therefore remain more specific. Paper 3 asks whether a \emph{scientifically scoped} structural object---with explicit QoI, context, directional role mapping, preserved and non-preserved properties, boundary conditions, evidence and negative-transfer history---can do two jobs with the same persistent identity: estimate cross-domain training redundancy and license or reject inference-time transfer. ReX supplies reusable latent experience; SkillGraph supplies dependency-aware retrieval; AgentGL supplies graph-conditioned learning plus graph-native inference. RAKL must either improve on these parents in the controlled cross-domain/semantic-decoy regime or narrow its claim to a different validity-governance niche.
